\documentclass[11pt,letterpaper]{article}

\usepackage[utf8]{inputenc}
\usepackage[T1]{fontenc}
\usepackage{lmodern}
\usepackage[margin=1.8cm]{geometry}
\usepackage{xcolor}
\usepackage{graphicx}
\usepackage{tabularx}
\usepackage{array}
\usepackage{hyperref}

\definecolor{navy}{HTML}{142A43}
\definecolor{teal}{HTML}{13806D}
\definecolor{gold}{HTML}{D99A1B}
\definecolor{mist}{HTML}{EDF4F5}
\definecolor{ink}{HTML}{263746}

\hypersetup{
  colorlinks=true,
  urlcolor=teal,
  linkcolor=teal,
  pdfauthor={Grupo 1 - CC3084 Data Science},
  pdftitle={Ficha de repositorio - Laboratorio 8 DuckDB}
}

\pagestyle{empty}

\setlength{\parindent}{0pt}
\setlength{\parskip}{6pt}

\newcommand{\labeltext}[1]{\textcolor{teal}{\bfseries\MakeUppercase{#1}}}

\begin{document}

\begin{minipage}[t]{0.68\textwidth}
  {\color{teal}\bfseries\large UNIVERSIDAD DEL VALLE DE GUATEMALA}\par
  {\color{ink}CC3084 - Data Science \textbar{} Sección 10 \textbar{} Grupo 1}\par
  {\color{ink}Segundo semestre 2026}
\end{minipage}%
\begin{minipage}[t]{0.32\textwidth}
  \raggedleft
  \colorbox{navy}{\parbox{0.76\linewidth}{\centering\color{white}\bfseries LABORATORIO 8\\[2pt]\normalsize DUCKDB}}
\end{minipage}

\vspace{8pt}
{\color{gold}\rule{\textwidth}{3pt}}

\vspace{10pt}
{\fontsize{25}{29}\selectfont\color{navy}\bfseries Análisis escalable de viajes NYC TLC}\par
{\large\color{teal}Yellow y Green Taxi - 2024, 2025 y corte parcial de 2026}

\vspace{10pt}
\colorbox{mist}{%
  \parbox{\dimexpr\textwidth-2\fboxsep\relax}{%
    \vspace{7pt}
    \labeltext{Fork oficial del equipo}\par
    \vspace{3pt}
    \url{https://github.com/DanielBarillasM/duckdb_Lab-8_Grupo-1_Sec-10}
    \vspace{7pt}
  }%
}

\vspace{12pt}
\begin{tabularx}{\textwidth}{>{\bfseries\color{navy}}X >{\raggedleft\arraybackslash}p{2.8cm}}
Integrante & Carné \\
Jorge Gabriel Palacios Sales & 231385 \\
Pablo Daniel Barillas Moreno & 22193 \\
Roberto Emiliano Otonie & 23968 \\
\end{tabularx}

\vspace{12pt}
\begin{minipage}[t]{0.48\textwidth}
  \labeltext{Contenido del repositorio}
  \begin{itemize}
    \item Descarga incremental de datos oficiales de NYC TLC.
    \item Doce consultas SQL directas sobre archivos Parquet.
    \item Notebook ejecutado con resultados y visualizaciones.
    \item Benchmark Parquet frente a tabla DuckDB.
    \item Siete indicadores y tablero reproducible en Metabase.
  \end{itemize}
\end{minipage}\hfill
\begin{minipage}[t]{0.48\textwidth}
  \labeltext{Reproducibilidad}
  \begin{itemize}
    \item Entorno definido mediante Docker Compose.
    \item Dependencias y transformaciones versionadas.
    \item Documentación completa en el README.
    \item Datos y credenciales excluidos del repositorio Git.
    \item Evidencia y mapa de cumplimiento de la rúbrica.
  \end{itemize}
\end{minipage}

\vfill
\colorbox{navy}{%
  \parbox{\dimexpr\textwidth-2\fboxsep\relax}{%
    \vspace{7pt}
    \centering\color{white}
    \textbf{Material de entrega}\par
    El enlace anterior conduce al fork versionado con el código, consultas, libreta, documentación, benchmark y evidencia del tablero.
    \vspace{7pt}
  }%
}

\end{document}
